\section{Background and Threat Model}
\label{sec:background}

\subsection{DER Operations and Standards}
\label{sec:bg-der}
A modern Distributed Energy Resource installation comprises a smart inverter
and its associated controller, a Battery Energy Storage System (BESS) where
present, and a Remote Terminal Unit (RTU) or Intelligent Electronic Device
(IED) bridging field-area communications to the upstream DER Management
System (DERMS) and Energy Management System (EMS). The control surface is
defined by IEEE Std~1547-2018, which specifies normal operating ranges for
voltage and frequency, ride-through requirements, and ramp-rate limits, and
by IEC~61850 (GOOSE/MMS) and DNP3 outstation profiles, which carry
telemetry and dispatch commands at the protocol level; IEEE Std
1547.3-2023 extends this surface with end-to-end DER cybersecurity
guidance~\cite{IEEE15473_2023}, and NIST SP 800-82r3 states the
OT-specific constraints (availability first, bounded latency, physical
consequence) under which any response mechanism must
operate~\cite{NIST80082r3}. Because each surface implies a different
attacker capability, our threat model and our scenario library both key
off these standards rather than off generic IT taxonomies.

\subsection{Threat Model}
\label{sec:bg-threat}
\textbf{Scope.} In scope are cyber attacks against the DER cyber-physical
control loop reachable from the utility-side network or from a compromised
aggregator account, including manipulated telemetry, forged commands, and
replayed sessions over IEC~61850 or DNP3. Out of scope are physical
hardware tampering, GPS-timing spoofing, pure ransomware impacting only
IT systems, and supply-chain firmware compromise of the firmware-attestation
pipeline (treated separately).

\textbf{Asset / surface matrix.} Five primary asset classes (smart inverter,
BESS controller, RTU/IED, DERMS/EMS, and the communications bus itself) are
exposed via different protocol surfaces; if owned, an attacker may
respectively manipulate reported $P/Q$ and trip thresholds, distort SOC
schedules, hijack telemetry aggregation, push spoofed dispatch, or replay
and reorder messages.

\textbf{STRIDE \texttimes{} {ATT\&CK} for ICS.} We map each STRIDE category
to the most relevant MITRE ATT\&CK for ICS technique IDs against each asset
class. Spoofing maps primarily to \emph{Valid Accounts} (T0859) and
\emph{Spoof Reporting Message} (T0856); Tampering to \emph{Manipulation of
View} (T0832) and \emph{Manipulation of Control} (T0831); Denial of Service
to T0814; and Elevation of Privilege to \emph{Exploit Public-Facing
Application} (T0890). The condensed mapping table is reproduced in
\texttt{docs/threat\_model.md} and feeds the attack-class taxonomy used
throughout this paper, namely \{\textsf{none}, \textsf{fdi}, \textsf{replay},
\textsf{command\_spoof}, \textsf{dos},
\textsf{firmware}\}~\cite{Zografopoulos2023_DER_CybersecurityOutlook,Hasan2024_ML_SG_CPS,Aslam2025_AI_ICS}.

\subsection{Trust Boundary and Assumptions}
\label{sec:bg-trust}
We state the trust boundary explicitly because it determines what the
architecture can and cannot claim.
\emph{Trusted computing base (implemented deterministically and tested):} the
FeatureView, the Evidence Gate, the fixed action registry, the safe-action
projection,
the irreversible-action escalation rule, the deadline fallback, and the
audit chain. \emph{Untrusted:} the LLM (and its adapter, prompt, memory,
and stated confidence), and any upstream IDS that supplies alerts.
\emph{Attacker knowledge/access:} the attacker may craft telemetry and
protocol frames reaching the detector, inject text into any field the LLM
reads (log lines, alert bodies, memory), and knows the architecture; the
attacker cannot alter the deterministic layer's code or the defender's
grid model. \emph{Evidence assurance is layered, and we state the layers
explicitly rather than conflating them.} Every gate feature is (i)
\emph{independent of LLM output} --- computed deterministically, so a
manipulated model cannot move it. It is (ii) \emph{wire-observable} ---
derived from packets, timestamps, sequence numbers, and reported power, or
from the defender's rated-capacity model. Observability is \emph{not}
authentication: on an unauthenticated bus an attacker who can inject or
modify network traffic can forge much of this evidence (a spoofed command
packet, a plausible reported-power series), so wire-observable evidence is
robust to a \emph{compromised model} but not to a \emph{compromised
network}. Only where a MAC, signature, or redundant physical measurement
is present is the evidence (iii) \emph{authenticated}; we implement no
cryptographic check and therefore claim only (i) and, for the
rating-plausibility residual, a physics cross-check that a purely
message-layer attacker cannot satisfy. We model the upstream IDS as an
\emph{imperfect} component and evaluate the Gate under detector
false-positive, false-negative, and latency perturbations. \emph{Human boundary:} HITL approval transfers authority to
a human operator; ``zero autonomous irreversible execution'' means the
system never executes an irreversible action \emph{without} human
approval --- it does \emph{not} guarantee that a human-approved action is
safe. \emph{Unsafe} means an irreversible or class-inappropriate physical
action executed autonomously, or an out-of-registry action.
\emph{Out of scope:} physical hardware tampering, cryptographic key
compromise, an adaptive white-box attacker that co-optimizes against the
deterministic layer, and a stealthy attacker able to forge internally
consistent protocol evidence (sequence numbers, timestamps, and physical
response all mutually consistent).

\subsection{Runtime-Assurance Objective}
\label{sec:bg-objective}
Given this threat model, the system objective is deliberately not
``classify attacks accurately.'' It is: (i)~no unsafe action may reach
the feeder even when the advisory model is wrong or adversarially
steered; (ii)~unnecessary mitigation on benign operating days must be bounded and
kept low (the measured false-open rate is nonzero, not zero); (iii)~a
bounded response must be available within OT deadlines regardless of
model latency; and (iv)~every decision must be auditable. The architecture is
designed so that autonomous policy containment does not depend on model
classification accuracy; service and reliability outcomes still depend on
evidence availability and gate detection. \emph{Safety in this study denotes
action-policy containment within the stated registry and threat model, not formal
invariance of feeder physical states.} We therefore distinguish throughout:
action-policy containment (what may execute), cyber-physical performance (ENS,
curtailment, voltage), detection coverage (gate behaviour), and formal safety
guarantees (which we do not claim).

\subsection{Physical-Impact Metrics}
\label{sec:bg-metrics}
We score every detector-error against a vector of physical-impact
metrics, all reported in their native units. We deliberately do
\emph{not} aggregate them into a monetary score: monetary weights vary
by jurisdiction, by tariff structure, and by the operator's posture,
and any single weighting gives the impression of a precision the
underlying weights do not have. Each component metric is a pure
function over the time-series emitted by the harness; the energy
values are computed by trapezoidal integration with units of kWh.

\begin{itemize}
  \item Energy Not Supplied (ENS) [kWh] --- integral of unmet demand
    over the run window; the headline reliability metric.
  \item Curtailed energy ($E_{\mathrm{curt}}$) [kWh] --- integral of
    available-but-not-dispatched generation; the headline
    over-mitigation metric.
  \item Voltage band excursion ($\phi_{\mathrm{volt}}$) [fraction of
    (bus,\,time) samples outside ANSI~C84.1].
  \item Frequency deviation $\max|\Delta f|$ [Hz] --- the safety-margin
    metric against IEEE~1547 ride-through bounds.
  \item IEEE~1547 ramp-rate violations [count].
\end{itemize}
Detection metrics (per-class precision/recall/F1, macro-F1) are
reported alongside the physical metrics and have a precise, separate
interpretation: they describe \emph{whether} an attack was correctly
classified, while the physical metrics describe the
\emph{consequence} of the action that followed. Sec.~\ref{sec:eval}
reports both, with paired bootstrap confidence intervals on the
ENS-headline metric.
